\chapter{不同牌力的处理框架}

\section{本章导入}

本章从单手牌思维上升到牌力分类框架。扑克是范围的游戏，同一手牌在不同牌面、位置和对手行动下可能承担不同功能。初学者需要先学会给自己的手牌分类，再决定它适合下注、过牌、跟注、加注还是弃牌。

本章将为后续下注目的、激进策略和弃牌跟注打下基础。只有理解不同牌力在整体范围中的角色，才能避免把所有牌都当成“强或弱”这种过于粗糙的二分法。

\section{核心概念}

本章将介绍坚果牌、强价值牌、中等摊牌价值、弱摊牌价值、强听牌、弱听牌和空气牌。读者需要理解，牌力分类不是固定标签，而是随公共牌和行动线变化的动态判断。

\section{新手常见误区}

新手常把顶对自动当成强牌，把听牌自动当成必须追的牌，也常低估中等摊牌价值的过牌意义。本节后续将说明不同牌力的主要任务。

\section{实战思考框架}

每一街行动前，可以先问：我的牌属于价值、摊牌价值、听牌还是空气牌？它需要保护吗？它能从更差牌获得价值吗？它能让更好牌弃牌吗？

\section{牌例拆解}

\subsection{牌例一：待补充}

本牌例将用于展示同一手顶对在不同牌面结构下的牌力变化。

\subsection{牌例二：待补充}

本牌例将用于展示强听牌和弱听牌在进攻策略上的差异。

\section{本章总结}

牌力分类是框架化决策的重要工具。读者应学会先识别手牌功能，再选择对应行动，而不是直接根据情绪决定进退。

\section{课后训练}

请选取五个翻牌局面，把自己的手牌分别归入本章的牌力类别，并写出最自然的行动方案。
